\documentclass[10pt,a4paper]{amsart}
\usepackage[margin=19mm]{geometry}
\usepackage{amsmath,amssymb,mathtools}
\usepackage[scr=rsfs,cal=euler]{mathalfa}
\usepackage{xfrac}
\usepackage[unicode,hidelinks,bookmarks=false]{hyperref}
\newcommand{\ie}{i.e.\ }
\newcommand{\cf}{cf.\ }
\newcommand{\resp}{respectively\ }
\newcommand{\Subsection}{\subsection}
\providecommand{\nobreakdash}{\nobreak\hbox{-}}
\setlength{\emergencystretch}{2em}
\multlinegap0pt
\usepackage[shortlabels]{enumitem}
\usepackage{amssymb}
\usepackage{tikz}
\usepackage{enumerate}
\usepackage{bm}
\usepackage{braket}
\usepackage{xcolor}
\usepackage{stackengine}
\newtheorem{definition}{Definition}
\newtheorem{theorem}{Theorem}
\newtheorem{proposition}{Proposition}
\newtheorem{remark}{Remark}
\newtheorem{lemma}{Lemma}
\newcommand{\El}{\underline{E}}
\newcommand{\Ga}{\mathit{\Gamma}}
\newcommand{\Gu}{\overline{\Ga}}
\newcommand{\TVinText}{\frac12 \big\|\tPy - P_Y^n \big\|_1}
\newcommand{\carY}{|\mY|}
\newcommand{\iQ}{\iota(Q_{XY})}
\newcommand{\mB}{\mathcal{B}}
\newcommand{\mC}{\mathcal{C}}
\newcommand{\mG}{\mathcal{G}}
\newcommand{\mP}{\mathcal{P}}
\newcommand{\mT}{\mathcal{T}}
\newcommand{\mX}{\mathcal{X}}
\newcommand{\mY}{\mathcal{Y}}
\newcommand{\nl}{\mathrm{nl}}
\newcommand{\su}{\mathrm{uni}}
\newcommand{\tP}{\tilde{P}}
\newcommand{\tPy}{\tP_{Y^n|\mC}}
\newcommand{\tPyQ}{\tP_{Y^n|\mC\sim q}}
\title{On the Strong Converse Exponent and Error Exponent of the Classical Soft Covering\vspace{2ex}}
\author{Xingyi He, S. Sandeep Pradhan,, Andreas Winter}
\date{Source: arXiv:2409.18307v4 (CC BY 4.0)}
\begin{document}
\maketitle

\noindent\emph{Source and licence.} On the Strong Converse Exponent and Error Exponent of the Classical Soft Covering\vspace{2ex}. Version arXiv:2409.18307v4 (CC BY 4.0), distributed by arXiv under the Creative Commons Attribution 4.0 International licence, http://creativecommons.org/licenses/by/4.0/, as recorded in the arXiv licence metadata for this exact version and shown on the abstract page https://arxiv.org/abs/2409.18307v4. Full licence notice: \texttt{licenses/arxiv-2409.18307-CC-BY-4.0.txt}.\par\smallskip

\noindent\textit{Editorial scope.} Counted unit: the article's theorem thm-nl-El (source0.tex lines 627--634). The article's complete original proof of it is delivered with it (source0.tex lines 1280--1391, one \texttt{proof} environment of 112 lines, which the article places away from the statement). No mathematics is written, repaired or re-derived: the statement and the proof are the article's own lines, in the article's own notation, and no retained line is substituted or re-flowed.  Retained uncounted as the source-local context the proof needs: lines 342--350; lines 1027--1030; lines 1182--1188; lines 1203--1206; lines 1269--1275; lines 1813--1816.  Nothing was omitted from any retained span, so the package carries no omission record.  No retained line contains a citation, so the wrapper carries no bibliography and no bibliography entry.  No retained line contains a figure environment, an image-inclusion command or drawing code, so no image is delivered and the package declares no asset dependency.  Editorial formatting only, in addition: an AMS \texttt{amsart} preamble that restates the article's own declarations and macros used by the retained lines, namely the environments \texttt{definition}, \texttt{theorem}, \texttt{proposition}, \texttt{remark}, \texttt{lemma} on separate counters, together with the standard packages those lines need.  The retained environments print in this document as Definition 1, Theorem 1, Proposition 1, Remark 1, Lemma 1 in order of appearance, whereas the article prints the same items with the numbering of its own full document.  The complete native source of the article as distributed in the arXiv e-print for this exact version is bundled unchanged as \texttt{sources/165897/source0.tex}.
\begin{definition}[Expectation of the information density]\label{def-iQ}
    Given $Q_{XY}\in\mP(\mX\mY)$, the expectation of the information density under $Q_{XY}$ is defined as
    $$ \iota(Q_{XY}) := \mathbb{E}_{Q_{XY}} [\iota_{X;Y}] = 
    \begin{cases}
        \displaystyle{ \sum_{x,y} Q_{XY}(x,y) \log\frac{W_{Y|X}(y|x)}{P_Y(y)} } 
            & \text{if } V\ll W \\
        -\infty & \text{otherwise}
    \end{cases}. $$
\end{definition}

\begin{proposition}[Achievability]\label{prop-sc-achi} The strong converse exponent  $\Ga_\su(R)$ in the uniform formulation can be bounded from above as follows:
    $$ \Ga_\su(R) \leq \Gu(R) := \min_{Q_{XY}\in\mP(\mX\mY): I(Q_X;V) \leq R} \left[D(Q_Y\|P_Y) + \big| R - \iQ \big|^+ \right], $$ 
    where $\iQ$ is defined in Definition \ref{def-iQ}.
\end{proposition}

\begin{remark}[Choice of alphabet]\label{rmk-chooseX} 
% In the proof of achievability, we choose some subset of $\mX$ and construct a code accordingly. In the proof of converse, we are always given a code. 
 If $w(\cdot)$ is a bijection, then, without loss of generality, we assume that  $\mX=\mY$.
     If $w(\cdot)$ is a surjection, for each $y\in\mY$, in the achievability proof we may select a single representative $x\in w^{-1}(y)$ and construct an alphabet $\mX_\mC \subset\mX$ consisting of these representatives and hence satisfying $|\mX_\mC| = |\mY|$. In the converse proof, given any code $\mC$, we first proceed with a restriction of symbols: if $y_1 = w(x_1) = w(x_2)$, we replace every occurrence of $x_2$ in the codewords with $x_1$. Since both $x_1$ and $x_2$ map to the same output symbol $y_1$, this modification leaves the output sequences unchanged and, consequently, does not affect the induced distribution $\tPyQ$ or $\tPy$. Under this restriction, the code alphabet set reduces to some $\mX_\mC \subset\mX$ such that $|\mX_\mC|=|\mY|$ and $w(\cdot)$ can be regarded bijective. 
Thus, in both cases, one can always identify an input alphabet set $\mX_\mC$ such that $|\mX_\mC|= |\mY|$ for encoding. 
% With this choice, $w(\cdot)$ can be treated as a bijection, therefore simplifying our discussion.
\end{remark}

\begin{equation}\label{badset}
    \mB:= \left\{y^n\in\mY^n: \frac{1}{M} \geq 2 P_Y^n(y^n) \right\} 
    = \bigsqcup_{Q_Y\in\mP_n(\mY): D(Q_Y\|P_Y) + H(Q_Y) \geq R + \frac1n} \mT_{Q_Y}.
\end{equation}

\begin{lemma}\label{lem-N1N2}
    When $R\geq H(P_Y)$, we have $N_1(R) \geq N_2(R)$,  where
    \begin{align*}
        N_1(R) &= \max_{Q_Y\in\mP(\mY): D(Q_Y\|P_Y) + H(Q_Y)\leq R} H(Q_Y), \\
        N_2(R) &= \max_{Q_Y\in\mP(\mY): D(Q_Y\|P_Y) + H(Q_Y)\geq R} \left[R-D(Q_Y\|P_Y)\right].
    \end{align*}
\end{lemma}

\begin{equation}\label{cuff-minQy}
    \min_{Q_Y\in\mP(\mY)} \left[ D(Q_Y\|P_Y) - \sum_y Q_Y(y) F(y) \right]
    = -\log\left( \sum_y P_Y(y) \ 2^{F(y)} \right)
\end{equation}

\begin{theorem}[Achievability for noiseless channels]\label{thm-nl-El} For noiseless channels under the uniform formulation, we have the following lower bound on $E_\su^\nl(R)$.
    \begin{align*}
        E_\su^\nl(R) \geq \El^\nl(R)
        := \ & \min_{Q_Y\in\mP(\mY)} 
        \left\{ D(Q_Y\|P_Y) + \big| R - D(Q_Y\|P_Y) - H(Q_Y) \big|^+ \right\} \\
        = \ & \max_{\alpha\geq1} \left\{ \frac{\alpha-1}{\alpha} \left[R - H_{\frac{1}{\alpha}}(P_Y)\right] \right\},
    \end{align*}
\end{theorem}

\begin{proof}[Proof of Theorem \ref{thm-nl-El}]
Taking inspiration from the proof of Proposition \ref{prop-sc-achi} regarding the achievability of the strong converse exponent, we provide a deterministic code construction for the problem at hand. 
Let us first assume that $R\geq H(P_Y)$. Define
\begin{equation}\label{goodset}
    \mG := \left\{ y^n\in\mY^n: MP_Y^n(y^n)\geq1 \right\}
    = \bigsqcup_{Q_Y\in\mP_n(\mY): D(Q_Y\|P_Y) + H(Q_Y)\leq R} \mT_{Q_Y}.
\end{equation}
$\mG$ is actually a `good' set in the sense that all sequences in $\mG$ satisfy that $\lfloor MP_Y^n(y^n)\rfloor\geq1$. Thus, it is always preferable to cover these sequences, since doing so will yield nonzero $k(y^n)$ values that can align with $MP_Y^n(y^n)$. In fact, ignoring the factor of two, $\mG$ is approximately the complement of the `bad' set $\mB$ defined in \eqref{badset}. Take any $\epsilon>0$. The size of $\mG$ is bounded by
\begin{align*}
    |\mG| &= \sum_{Q_Y\in\mP_n(\mY): D(Q_Y\|P_Y) + H(Q_Y)\leq R} \mT_{Q_Y} 
    \geq \max_{Q_Y\in\mP_n(\mY): D(Q_Y\|P_Y) + H(Q_Y)\leq R} \mT_{Q_Y} \\
    &\overset{a}{\geq} (n+1)^{-\carY} \max_{Q_Y\in\mP_n(\mY): D(Q_Y\|P_Y) + H(Q_Y)\leq R} 2^{nH(Q_Y)} \\
    % &= (n+1)^{-\carY} \exp_2 \left\{ n \max_{Q_Y\in\mP_n(\mY): D(Q_Y\|P_Y) + H(Q_Y)\leq R} H(Q_Y) \right\} \\
    &\overset{b}{\geq} (n+1)^{-\carY} \ 2^{n \left[N_1(R)-\epsilon\right]},
\end{align*}
where $(a)$ follows from property of types. In $(b)$, the maximization is over all types, so it differs from the maximization over distributions (namely $N_1(R)$ as defined in Lemma \ref{lem-N1N2}) by at most $\epsilon$ for sufficiently large $n$.

To design a good code, we adopt the alphabet choice described in Remark \ref{rmk-chooseX}. With this choice, $k(y^n)$ represents the number of codewords mapped one-to-one to $y^n$, and hence can be directly constructed via a repetition of codewords. Our goal is to judiciously design $k(y^n)$  so that it can satisfy $\lfloor MP_Y^n(y^n) \rfloor \leq k(y^n) \leq \lceil MP_Y^n(y^n) \rceil + \mathrm{poly}(n)$ under the constraint that $\sum_{y^n} k(y^n) = M$. Namely, let us first consider the following two codes:
\begin{align*}
    \mC_1: \quad 
    & k(y^n) = \lfloor MP_Y^n(y^n) \rfloor \quad \forall y^n\in\mY^n, \quad
    M_1 := |\mC_1| = \sum_{y^n} \ \lfloor MP_Y^n(y^n) \rfloor. \\
    \mC_2: \quad 
    & k(y^n) = 
    \begin{cases}
        \lceil MP_Y^n(y^n) \rceil & \text{if } y^n\in\mG \\
        0 & \text{if } y^n\notin\mG
    \end{cases}, \quad
    M_2 := |\mC_2| = \sum_{y^n\in\mG} \lceil MP_Y^n(y^n) \rceil.
\end{align*}
Clearly $M_1\leq M$, meaning that we can use $\mC_1$ to cover all sequences $y^n\in\mY^n$ and there will be some codewords left. Note that $\mG^c := \mY^n\backslash\mG$ is not covered since $\lfloor MP_Y^n(y^n) \rfloor = 0$ for all $y^n\in\mG^c$.
If $M_2\geq M$, then after applying $\mC_1$, each sequence in $\mG$ can be covered at most once more. Hence, there exists a code $\mC_3$ as follows, with $|\mC_3| = M$:
\begin{align*}
    \mC_3: \quad 
    & k(y^n) = 
    \begin{cases}
        \lfloor MP_Y^n(y^n) \rfloor \text{ or } \lceil MP_Y^n(y^n) \rceil  & \text{if } y^n\in\mG \\
        0 & \text{if } y^n\notin\mG
    \end{cases}.
\end{align*}
If $M_2<M$, then after applying $\mC_1$, each sequence in $\mG$ can still be covered more than once. Equivalently, we use $\mC_2$ to cover $\mY^n$, and there will be some codewords left. The number of those codewords that are left is
\begin{align*}
    \Delta M &:= M - M_2 
    = M - \sum_{y^n\in\mG} \lceil MP_Y^n(y^n) \rceil 
    \leq M - \sum_{y^n\in\mG} MP_Y^n(y^n) 
    = M P_Y^n(\mG^c) \\
    &= \sum_{Q_Y\in\mP_n(\mX): D(Q_Y\|P_Y) + H(Q_Y)>R} MP_Y^n(\mT_{Q_Y}) \\
    &\overset{a}{\leq} \sum_{Q_Y\in\mP_n(\mY): D(Q_Y\|P_Y) + H(Q_Y)>R} 2^{n\left[R-D(Q_Y\|P_Y)\right]} \\
    &\overset{b}{\leq} (n+1)^{\carY} \max_{Q_Y\in\mP_n(\mY): D(Q_Y\|P_Y) + H(Q_Y)>R} 2^{n\left[R-D(Q_Y\|P_Y)\right]} \\
    &= (n+1)^{\carY} \exp_2 \left\{ n \max_{Q_Y\in\mP_n(\mY): D(Q_Y\|P_Y) + H(Q_Y)>R} \left[R-D(Q_Y\|P_Y)\right] \right\} \\
    &\leq (n+1)^{\carY} \ 2^{n N_2(R)},
\end{align*}
where $(a)$ and $(b)$ follow from properties of types, and $N_2(R)$ is defined in Lemma \ref{lem-N1N2}. 
Now let us consider using these $\Delta M$ codewords to continue covering $\mG$ after $\mC_2$ is applied. We can cover each sequence in $\mG$ by $\frac{\Delta M}{|\mG|}$ more times, which is at most
$$ \frac{\Delta M}{|\mG|} \leq (n+1)^{2\carY} 2^{n \left[N_2(R)-N_1(R) + \epsilon\right]} \leq (n+1)^{2\carY} 2^{n\epsilon} \leq 2^{2n\epsilon} $$
for all sufficiently large $n$. Here we have used $N_1(R)\geq N_2(R)$ when $R\geq H(P_Y)$ from Lemma \ref{lem-N1N2}. Hence, there exists a code $\mC_4$ as follows, with $|\mC_4| = M$:
\begin{align*}
    \mC_4: \quad 
    &
    \begin{cases}
        \lceil MP_Y^n(y^n) \rceil \leq k(y^n) \leq \lceil MP_Y^n(y^n) \rceil + 2^{2n\epsilon} & \text{if } y^n\in\mG \\
        k(y^n) = 0 & \text{if } y^n\notin\mG
    \end{cases}.
\end{align*}

According to the above discussions, either $\mC_3$ or $\mC_4$ is possible and has exactly $M$ codewords. Therefore, we can summarize that for all $\epsilon>0$ and sufficiently large $n$, there exists a code $\mC$ as follows, with $|\mC| = \sum_{y^n} k(y^n) = M$:
\begin{align*}
    \mC: \quad 
    &
    \begin{cases}
        \lfloor MP_Y^n(y^n) \rfloor \leq k(y^n) \leq \lceil MP_Y^n(y^n) \rceil + 2^{2n\epsilon} & \text{if } y^n\in\mG \\
        k(y^n) = 0 & \text{if } y^n\notin\mG
    \end{cases}.
\end{align*}
Choose $\mC$ to be our covering code, which yields that
\[
  \left| \frac{k(y^n)}{M} - P_Y^n(y^n) \right| 
   \leq \left\{\begin{array}{@{}l l@{}}
          \dfrac{1+2^{2n\epsilon}}{M} & \text{if } y^n \in \mG \\
          P_Y^n(y^n) & \text{if } y^n \notin \mG
        \end{array}\right\}
   \leq 2^{3n\epsilon} \min\left\{\frac{1}{M}, P_Y^n(y^n)\right\}. 
\]
Consequently, we have 
\begin{align*}
    \frac12 \left\|\tPy - P_Y^n\right\|_1
    &\leq 2^{3n\epsilon-1}  \sum_{y^n} \min\left\{\frac{1}{M}, P_Y^n(y^n)\right\} \\
    &= 2^{3n\epsilon-1} \sum_{Q_Y\in\mP_n(\mY)} |\mT_{Q_Y}| 
        \min\left\{ 2^{-nR}, 2^{-n\left[D(Q_Y\|P_Y) + H(Q_Y)\right]} \right\} \\
    &\overset{a}{\leq} 2^{3n\epsilon-1} \sum_{Q_Y\in\mP_n(\mY)} 2^{nH(Q_Y)} 
        \min\left\{ 2^{-nR}, 2^{-n\left[D(Q_Y\|P_Y) + H(Q_Y)\right]} \right\} \\
    &\overset{b}{\leq} \frac12(n+1)^{\carY} \ 2^{3n\epsilon} \max_{Q_Y\in\mP_n(\mY)} \min\left\{ 2^{-n[R - H(Q_Y)]}, 2^{-nD(Q_Y\|P_Y)} \right\} \\
    &= \frac12(n+1)^{\carY} \exp_2 \left\{-n \min_{Q_Y\in\mP_n(\mY)} \left[ D(Q_Y\|P_Y) + \big| R - D(Q_Y\|P_Y) - H(Q_Y) \big|^+ - 3\epsilon \right] \right\},
\end{align*}
where $(a)$ and $(b)$ follow from properties of types. The exponent is exactly $\El^\nl(R)$ as $\epsilon\to0$ and $n\to\infty$. 

Note that all the above discussions are based on the assumption that $R\geq H(P_Y)$. If $R<H(P_Y)$, we can just apply a trivial bound $\TVinText\leq1$ and the error exponent is $0$. Since $\El^\nl(R)$ also vanishes for $R<H(P_Y)$, it is thus an achievable error exponent for $R$ values both below and above $H(P_Y)$. Hence, the variational form in the theorem is proved.

To obtain the dual form, consider the following.
\begin{align*}
    \min_{Q_Y\in\mP(\mY)} &
        \left\{ D(Q_Y\|P_Y) + \big| R - D(Q_Y\|P_Y) - H(Q_Y) \big|^+ \right\} \\
    = \ & \max_{\lambda\in[0,1]} \ \min_{Q_Y\in\mP(\mY)} 
        \big\{ D(Q_Y\|P_Y) + \lambda\left[ R - D(Q_Y\|P_Y) - H(Q_Y) \right] \big\} \\
    = \ & \max_{\lambda\in[0,1]} \left\{\lambda R + \min_{Q_Y\in\mP(\mY)}
        \left[ D(Q_Y\|P_Y) + \lambda \sum_y Q_Y(y) \log P_Y(y) \right] \right\} \\
    \overset{a}{=} \ & \max_{\lambda\in[0,1]}
        \left\{ \lambda R - \log \sum_y P_Y^{1-\lambda}(y) \right\} \\
    \overset{b}{=} \ & \max_{\alpha\geq1} \left\{ \frac{\alpha-1}{\alpha} \left[R - H_{\frac{1}{\alpha}}(P_Y)\right] \right\},
\end{align*}
where $(a)$ follows from \eqref{cuff-minQy} and $(b)$ follows by setting $\alpha:= \frac{1}{1-\lambda}$.
\end{proof}

\end{document}
